\noindent\textbf{Supplementary Table S13E.} Exploratory Gevers medication follow-up. The first three analyses asked whether the pediatric Crohn portability boundary persisted after restricting the Gevers stool branch to samples explicitly flagged antibiotic-negative, mesalamine-negative, or steroid-negative. The next three analyses compared medication-positive with medication-negative Crohn-case samples among cases with known status for the same field. Biologics were not evaluable because no analyzed Gevers stool samples were flagged biologics-positive.

\noindent\textbf{Analysis:} Antibiotic-negative Gevers subset. The cohort was restricted to samples explicitly flagged antibiotic-negative, and Crohn disease was compared with healthy controls within that subset. The explicit antibiotic-negative subset contained 321 Crohn samples and 26 healthy-control samples.

\begin{center}
\scriptsize
\begin{tabularx}{\textwidth}{@{} l r r Y Y r r Y @{}}
\toprule
Mode & Depth & Positive/negative samples & Top dCST & Frequentist OR (95\% CI) & \(p\) & \(q\) & Bayesian OR (95\% CrI) \\
\midrule
Rebuilt & 1 & 206/26 & \emph{Faecalibacterium prausnitzii} & OR 2.76 [1.05, 8.1] & 0.036 & 0.107 & OR 2.74 [1.15, 7.19] \\
Rebuilt & 2 & 206/26 & \emph{Faecalibacterium prausnitzii}\dcstsep \emph{Dialister invisus} & OR 3.44 [0.802, 31.1] & 0.121 & 0.392 & OR 3.23 [0.944, 17.5] \\
Rebuilt & 3 & 206/26 & \emph{Blautia obeum}\dcstsep \emph{Bifidobacterium breve} & OR 0.285 [0.0843, 1.12] & 0.037 & 0.456 & OR 0.277 [0.0947, 0.913] \\
Rebuilt & 4 & 206/26 & \emph{Blautia obeum}\dcstsep \emph{Bifidobacterium breve} & OR 0.285 [0.0843, 1.12] & 0.037 & 0.456 & OR 0.28 [0.0967, 0.903] \\
AGP transfer & 1 & 206/26 & \emph{Segatella sp.} & OR 0.272 [0.057, 1.74] & 0.088 & 1.000 & OR 0.26 [0.0693, 1.18] \\
AGP transfer & 2 & 189/26 & \emph{Bacteroides dorei}\dcstsep \emph{Bifidobacterium breve} & OR 0 [0, 5.37] & 0.121 & 1.000 & OR 0.0261 [0.0000511, 0.785] \\
AGP transfer & 3 & 174/24 & \emph{Segatella sp.}\dcstsep \emph{Faecalibacterium prausnitzii}\dcstsep \emph{UCG-002 sp.} & OR 0.13 [0.00902, 1.88] & 0.073 & 1.000 & OR 0.128 [0.0171, 0.967] \\
AGP transfer & 4 & 173/24 & \emph{Segatella sp.}\dcstsep \emph{Faecalibacterium prausnitzii}\dcstsep \emph{UCG-002 sp.} & OR 0.131 [0.00908, 1.89] & 0.073 & 1.000 & OR 0.128 [0.0167, 0.999] \\
\bottomrule
\end{tabularx}
\end{center}
\bigskip

\noindent\textbf{Analysis:} Mesalamine-negative Gevers subset. The cohort was restricted to samples explicitly flagged mesalamine-negative, and Crohn disease was compared with healthy controls within that subset. The explicit mesalamine-negative subset contained 339 Crohn samples and 26 healthy-control samples.

\begin{center}
\scriptsize
\begin{tabularx}{\textwidth}{@{} l r r Y Y r r Y @{}}
\toprule
Mode & Depth & Positive/negative samples & Top dCST & Frequentist OR (95\% CI) & \(p\) & \(q\) & Bayesian OR (95\% CrI) \\
\midrule
Rebuilt & 1 & 224/26 & \emph{Faecalibacterium prausnitzii} & OR 2.75 [1.06, 8.07] & 0.036 & 0.109 & OR 2.73 [1.15, 7.11] \\
Rebuilt & 2 & 224/26 & \emph{Faecalibacterium prausnitzii}\dcstsep \emph{Dialister invisus} & OR 3.8 [0.893, 34.2] & 0.079 & 0.344 & OR 3.57 [1.05, 19.6] \\
Rebuilt & 3 & 224/26 & \emph{Blautia obeum}\dcstsep \emph{Bifidobacterium breve} & OR 0.282 [0.0847, 1.1] & 0.034 & 0.523 & OR 0.273 [0.0969, 0.883] \\
Rebuilt & 4 & 224/26 & \emph{Blautia obeum}\dcstsep \emph{Bifidobacterium breve} & OR 0.282 [0.0847, 1.1] & 0.034 & 0.523 & OR 0.275 [0.0951, 0.888] \\
AGP transfer & 1 & 224/26 & \emph{Segatella sp.} & OR 0.141 [0.0224, 1.02] & 0.026 & 0.715 & OR 0.137 [0.0307, 0.676] \\
AGP transfer & 2 & 198/26 & \emph{Segatella sp.}\dcstsep \emph{Faecalibacterium prausnitzii} & OR 0.187 [0.0204, 2.34] & 0.104 & 1.000 & OR 0.179 [0.0313, 1.24] \\
AGP transfer & 3 & 183/24 & \emph{Segatella sp.}\dcstsep \emph{Faecalibacterium prausnitzii}\dcstsep \emph{UCG-002 sp.} & OR 0.124 [0.00858, 1.79] & 0.067 & 1.000 & OR 0.121 [0.0162, 0.905] \\
AGP transfer & 4 & 182/24 & \emph{Segatella sp.}\dcstsep \emph{Faecalibacterium prausnitzii}\dcstsep \emph{UCG-002 sp.} & OR 0.125 [0.00863, 1.8] & 0.068 & 1.000 & OR 0.123 [0.0163, 0.916] \\
\bottomrule
\end{tabularx}
\end{center}
\bigskip

\noindent\textbf{Analysis:} Steroid-negative Gevers subset. The cohort was restricted to samples explicitly flagged steroid-negative, and Crohn disease was compared with healthy controls within that subset. The explicit steroid-negative subset contained 351 Crohn samples and 26 healthy-control samples.

\begin{center}
\scriptsize
\begin{tabularx}{\textwidth}{@{} l r r Y Y r r Y @{}}
\toprule
Mode & Depth & Positive/negative samples & Top dCST & Frequentist OR (95\% CI) & \(p\) & \(q\) & Bayesian OR (95\% CrI) \\
\midrule
Rebuilt & 1 & 236/26 & \emph{Faecalibacterium prausnitzii} & OR 2.85 [1.09, 8.32] & 0.022 & 0.067 & OR 2.83 [1.19, 7.32] \\
Rebuilt & 2 & 236/26 & \emph{Faecalibacterium prausnitzii}\dcstsep \emph{Dialister invisus} & OR 3.99 [0.941, 35.8] & 0.051 & 0.256 & OR 3.73 [1.1, 20.9] \\
Rebuilt & 3 & 236/26 & \emph{Blautia obeum}\dcstsep \emph{Bifidobacterium breve} & OR 0.287 [0.0874, 1.11] & 0.035 & 0.464 & OR 0.282 [0.0981, 0.915] \\
Rebuilt & 4 & 236/26 & \emph{Blautia obeum}\dcstsep \emph{Bifidobacterium breve} & OR 0.287 [0.0874, 1.11] & 0.035 & 0.464 & OR 0.28 [0.099, 0.897] \\
AGP transfer & 1 & 236/26 & \emph{Segatella sp.} & OR 0.168 [0.0304, 1.15] & 0.035 & 0.976 & OR 0.161 [0.0382, 0.8] \\
AGP transfer & 2 & 207/26 & \emph{Segatella sp.}\dcstsep \emph{Faecalibacterium prausnitzii} & OR 0.179 [0.0194, 2.24] & 0.097 & 1.000 & OR 0.175 [0.031, 1.19] \\
AGP transfer & 3 & 191/24 & \emph{Segatella sp.}\dcstsep \emph{Faecalibacterium prausnitzii}\dcstsep \emph{UCG-002 sp.} & OR 0.119 [0.00822, 1.71] & 0.062 & 1.000 & OR 0.118 [0.0157, 0.881] \\
AGP transfer & 4 & 190/24 & \emph{Segatella sp.}\dcstsep \emph{Faecalibacterium prausnitzii}\dcstsep \emph{UCG-002 sp.} & OR 0.119 [0.00826, 1.72] & 0.063 & 1.000 & OR 0.12 [0.0154, 0.926] \\
\bottomrule
\end{tabularx}
\end{center}
\bigskip

\noindent\textbf{Analysis:} Antibiotic-positive vs antibiotic-negative Crohn samples. Only Crohn-case samples with known antibiotic status were retained, and antibiotic-positive samples were compared with antibiotic-negative samples. Among Crohn-case samples with known antibiotic status, 49 were antibiotic-positive and 206 were antibiotic-negative.

\begin{center}
\scriptsize
\begin{tabularx}{\textwidth}{@{} l r r Y Y r r Y @{}}
\toprule
Mode & Depth & Positive/negative samples & Top dCST & Frequentist OR (95\% CI) & \(p\) & \(q\) & Bayesian OR (95\% CrI) \\
\midrule
Rebuilt & 1 & 49/206 & \emph{Blautia obeum} & OR 1.34 [0.625, 2.77] & 0.468 & 0.871 & OR 1.34 [0.666, 2.6] \\
Rebuilt & 2 & 49/206 & \emph{Blautia obeum} & OR 1.34 [0.625, 2.77] & 0.468 & 0.726 & OR 1.34 [0.675, 2.64] \\
Rebuilt & 3 & 49/206 & \emph{Faecalibacterium prausnitzii}\dcstsep \emph{Bacteroides vulgatus}\dcstsep \emph{Bacteroides fragilis} & OR 0 [0, 1.35] & 0.080 & 0.877 & OR 0.0668 [0.000131, 0.825] \\
Rebuilt & 4 & 49/206 & \emph{Faecalibacterium prausnitzii}\dcstsep \emph{Bacteroides vulgatus}\dcstsep \emph{Bacteroides fragilis} & OR 0 [0, 1.35] & 0.080 & 0.877 & OR 0.0694 [0.00013, 0.844] \\
AGP transfer & 1 & 49/206 & \emph{Enterococcus faecium} & OR 4.26 [1.6, 11.2] & 0.002 & 0.046 & OR 4.29 [1.77, 10] \\
AGP transfer & 2 & 34/189 & \emph{Enterococcus faecium} & OR 6.39 [2.31, 17.6] & 1.3e-04 & 0.005 & OR 6.49 [2.6, 16] \\
AGP transfer & 3 & 32/174 & \emph{Enterococcus faecium} & OR 6.4 [2.28, 17.9] & 1.5e-04 & 0.006 & OR 6.52 [2.6, 16.4] \\
AGP transfer & 4 & 32/173 & \emph{Enterococcus faecium} & OR 6.36 [2.27, 17.7] & 1.6e-04 & 0.006 & OR 6.41 [2.56, 16.2] \\
\bottomrule
\end{tabularx}
\end{center}
\bigskip

\noindent\textbf{Analysis:} Mesalamine-positive vs mesalamine-negative Crohn samples. Only Crohn-case samples with known mesalamine status were retained, and mesalamine-positive samples were compared with mesalamine-negative samples. Among Crohn-case samples with known mesalamine status, 31 were mesalamine-positive and 224 were mesalamine-negative.

\begin{center}
\scriptsize
\begin{tabularx}{\textwidth}{@{} l r r Y Y r r Y @{}}
\toprule
Mode & Depth & Positive/negative samples & Top dCST & Frequentist OR (95\% CI) & \(p\) & \(q\) & Bayesian OR (95\% CrI) \\
\midrule
Rebuilt & 1 & 31/224 & \emph{Blautia obeum} & OR 1.43 [0.564, 3.4] & 0.387 & 0.850 & OR 1.43 [0.618, 3.15] \\
Rebuilt & 2 & 31/224 & \emph{Faecalibacterium prausnitzii}\dcstsep \emph{Segatella sp.} & OR 3.51 [0.74, 13.7] & 0.058 & 0.292 & OR 3.6 [0.993, 11.6] \\
Rebuilt & 3 & 31/224 & \emph{Bacteroides vulgatus}\dcstsep \emph{[Ruminococcus] gnavus group gnavus} & OR 3.51 [0.74, 13.7] & 0.058 & 0.554 & OR 3.62 [0.989, 11.7] \\
Rebuilt & 4 & 31/224 & \emph{Bacteroides vulgatus}\dcstsep \emph{[Ruminococcus] gnavus group gnavus} & OR 3.51 [0.74, 13.7] & 0.058 & 0.554 & OR 3.61 [0.972, 11.8] \\
AGP transfer & 1 & 31/224 & \emph{Gemmiger sp.} & OR 6.41 [1.2, 31.8] & 0.015 & 0.423 & OR 6.5 [1.57, 25.1] \\
AGP transfer & 2 & 25/198 & \emph{Gemmiger sp.} & OR 7.23 [1.33, 36.6] & 0.011 & 0.224 & OR 7.43 [1.8, 29.6] \\
AGP transfer & 3 & 23/183 & \emph{Gemmiger sp.} & OR 7.36 [1.34, 37.6] & 0.010 & 0.232 & OR 7.56 [1.82, 29.8] \\
AGP transfer & 4 & 23/182 & \emph{Gemmiger sp.} & OR 7.32 [1.34, 37.3] & 0.011 & 0.231 & OR 7.42 [1.78, 30.1] \\
\bottomrule
\end{tabularx}
\end{center}
\bigskip

\noindent\textbf{Analysis:} Steroid-positive vs steroid-negative Crohn samples. Only Crohn-case samples with known steroid status were retained, and steroid-positive samples were compared with steroid-negative samples. Among Crohn-case samples with known steroid status, 19 were steroid-positive and 236 were steroid-negative.

\begin{center}
\scriptsize
\begin{tabularx}{\textwidth}{@{} l r r Y Y r r Y @{}}
\toprule
Mode & Depth & Positive/negative samples & Top dCST & Frequentist OR (95\% CI) & \(p\) & \(q\) & Bayesian OR (95\% CrI) \\
\midrule
Rebuilt & 1 & 19/236 & \emph{Faecalibacterium prausnitzii} & OR 0.556 [0.179, 1.59] & 0.244 & 0.419 & OR 0.556 [0.205, 1.41] \\
Rebuilt & 2 & 19/236 & \emph{Faecalibacterium prausnitzii}\dcstsep \emph{Dialister invisus} & OR 0 [0, 0.67] & 0.009 & 0.045 & OR 0.0357 [0.0000705, 0.436] \\
Rebuilt & 3 & 19/236 & \emph{Bacteroides vulgatus}\dcstsep \emph{[Ruminococcus] gnavus group gnavus} & OR 4.2 [0.676, 18.6] & 0.062 & 0.535 & OR 4.37 [0.969, 15.6] \\
Rebuilt & 4 & 19/236 & \emph{Bacteroides vulgatus}\dcstsep \emph{[Ruminococcus] gnavus group gnavus} & OR 4.2 [0.676, 18.6] & 0.062 & 0.535 & OR 4.33 [0.974, 15.7] \\
AGP transfer & 1 & 19/236 & \emph{Gemmiger sp.} & OR 7.07 [1.05, 37] & 0.022 & 0.640 & OR 7.47 [1.55, 30.7] \\
AGP transfer & 2 & 16/207 & \emph{Faecalibacterium prausnitzii}\dcstsep \emph{Bifidobacterium breve} & OR 6.46 [1.29, 27.1] & 0.012 & 0.373 & OR 6.7 [1.7, 23.2] \\
AGP transfer & 3 & 15/191 & \emph{Faecalibacterium prausnitzii}\dcstsep \emph{Bifidobacterium breve} & OR 6.47 [1.28, 27.5] & 0.012 & 0.415 & OR 6.69 [1.7, 24] \\
AGP transfer & 4 & 15/190 & \emph{Faecalibacterium prausnitzii}\dcstsep \emph{Bifidobacterium breve} & OR 6.43 [1.27, 27.4] & 0.012 & 0.411 & OR 6.69 [1.71, 23.9] \\
\bottomrule
\end{tabularx}
\end{center}
\bigskip

